%%%%%%%%%%  scan idx 190  |  folio 177  %%%%%%%%%%
D'%
autre part, si $A$ est mou, tout $A$-Module est mou (T.F II 3.7.1). Comme exemples de faisceaux d'anneaux mous, citons le faisceau des fonctions continues (réelles ou complexes) sur un espace paracompact, le faisceau des fonctions $C^r$ sur une variété de classe $C^r$.

\textbf{Proposition 2.3.2.} Soit $S$ un espace compact, muni d'un faisceau d'anneaux mou $A$. On note $\Mod(S)$ la catégorie des $A$-Modules à gauche, $\mathrm D(S)=\mathrm D(\Mod(S))$.
\begin{enumerate}[label=\alph*)]
\item La catégorie $\Lib(S)$ (resp. $\Loclib(S)$) (2.0) est relevable dans $\Mod(S)$ et quasi-stable (resp. stable) par noyau d'épimorphisme (I 1.2)\footnote{Cette assertion vaut plus généralement pour $S$ paracompact.}.
\item Soient $E\in\Ob\mathrm D(S)$ et $n\in\mathbb Z$. Pour que $E$ soit $n$-pseudo-cohérent (resp. pseudo-cohérent) il faut et il suffit que $E$ le soit relativement à $\Lib(S)$ ou à $\Loclib(S)$. Les foncteurs canoniques
\[
\mathrm K^{-}(\Lib(S))\longrightarrow \mathrm D(S)
\]
et
\[
\mathrm K^{-}(\Loclib(S))\longrightarrow \mathrm D(S)
\]
sont pleinement fidèles et ont pour image essentielle $\mathrm D^{-}(S)_{\coh}$ (I 2.15).
\item Soient $E\in\Ob\mathrm D(S)$, et $[m,n]$ un intervalle de $\mathbb Z$. Pour que $E$ soit d'amplitude parfaite contenue dans $[m,n]$ (4.8) il faut et il suffit qu'il existe un quasi-isomorphisme $F\to E$, où $F$ est un complexe tel que $F^i=0$ pour $i\notin[m,n]$, $F^i\in\Ob\Lib(S)$ pour $i>m$ et $F^m\in\Ob\Loclib(S)$. Le foncteur canonique
\[
\mathrm K^b(\Loclib(S))\longrightarrow\mathrm D(S)
\]
est pleinement fidèle et a pour image essentielle $\mathrm D(S)_{\parf}$ (4.9).
\item Notons $\Mod(A(S))$ la catégorie des $A(S)$-modules à gauche, $\mathrm D(A(S))$ la catégorie dérivée correspondante. Le foncteur
\[
H^0(S,.):\Mod(S)\longrightarrow\Mod(A(S))
\]
est exact et une équivalence de catégories. En outre, il
définit des 
